\documentclass[aspectratio=169,11pt]{beamer}
 
% ---------------------------------------------------------------- thema
\usetheme[
  progressbar=frametitle,
  numbering=fraction,
  sectionpage=progressbar,
  block=fill
]{metropolis}
 
\usepackage[T1]{fontenc}
\IfFileExists{dutch.ldf}{%
  \usepackage[dutch]{babel}%
  \AtBeginDocument{\shorthandoff{"}}
}{}
 
\IfFileExists{plex-sans.sty}{\usepackage[sfdefault]{plex-sans}}{}
\IfFileExists{plex-mono.sty}{\usepackage[scale=0.95]{plex-mono}}{}
 
\usepackage{tikz}
\usetikzlibrary{positioning,arrows.meta,calc}
\usepackage{tikz-cd}
\usepackage{listings}
\usepackage{booktabs}
\usepackage{amssymb}
 
\definecolor{kulblauw}{HTML}{1D8DB0}
\definecolor{kuldonker}{HTML}{00407A}
 
\setbeamercolor{frametitle}{bg=kuldonker, fg=white}
\setbeamercolor{progress bar}{fg=kulblauw, bg=kulblauw!20}
\setbeamercolor{alerted text}{fg=kulblauw}
\setbeamercolor{standout}{bg=kuldonker, fg=white}
\setbeamercolor{block title}{bg=kulblauw!20, fg=kuldonker}
\setbeamercolor{block body}{bg=kulblauw!6}
\setbeamercolor{block title alerted}{bg=orange!25, fg=black}
\setbeamercolor{block body alerted}{bg=orange!8}
 
% ---------------------------------------------------------------- code
\definecolor{codebg}{HTML}{EEF4F8}
\definecolor{codecomment}{HTML}{6A7A86}
\definecolor{codestring}{HTML}{2E7D32}
\definecolor{codetype}{HTML}{C2185B}   % Int, Char, ...
\definecolor{codefunc}{HTML}{4A6FD1}   % replicate, swap, ...
\definecolor{codeop}{HTML}{E5395C}     % :: -> => <-
\definecolor{codenum}{HTML}{D9730D}    % 3, 10, ...
\definecolor{codekw}{HTML}{8E44AD}     % where, let, ...
 
\newcommand{\hsop}[1]{{\color{codeop}#1}}
\newcommand{\hsnum}[1]{{\color{codenum}#1}}
 
\lstdefinestyle{hs}{
  language=Haskell,
  sensitive=true,
  % keywords
  keywords={where,let,in,import,if,then,else,case,of,data,type,module,
            newtype,class,instance,deriving,do},
  keywordstyle=\color{codekw}\bfseries,
  % types en typeclasses
  keywords=[2]{Int,Char,String,Bool,Integer,Double,Float,Maybe,Either,
               Ord,Eq,Show,Num,True,False,Just,Nothing,Left,Right,IO,
               Functor,Applicative,Monad,Semigroup,Monoid,Foldable,
               Traversable,Bounded,Sum,Product,Max,Min,Any,All,First,
               NonEmpty,Identity,Const,State,Ordering},
  keywordstyle=[2]\color{codetype},
  % Prelude / user-defined functions
  keywords=[3]{replicate,flip,map,filter,length,zip,unzip,zipWith,span,
               break,sort,sortOn,curry,uncurry,swap,optellen,fst,snd,head,
               tail,sum,product,lines,words,unwords,read,show,foldr,foldl,
               reverse,concat,concatMap,take,drop,takeWhile,dropWhile,
               elem,all,any,maximum,minimum,abs,id,const,not,
               fmap,pure,traverse,sequenceA,foldMap,mempty,mconcat,sconcat,
               mapM,mapM_,join,return,get,put,modify,runState,evalState,
               execState,getSum,getProduct,getMax,getAny,liftA2,toList,
               null,even,div,putStrLn,getLine,readFile,maybeToList,
               liftM,ap,fold,unit,pair},
  keywordstyle=[3]\color{codefunc},
  basicstyle=\ttfamily\footnotesize,
  commentstyle=\color{codecomment}\itshape,
  stringstyle=\color{codestring},
  backgroundcolor=\color{codebg},
  columns=fullflexible, keepspaces=true,
  showstringspaces=false, upquote=true,
  xleftmargin=6pt, xrightmargin=6pt,
  framexleftmargin=6pt, framexrightmargin=6pt,
  framextopmargin=3pt, framexbottommargin=3pt,
  frame=single, rulecolor=\color{codebg},
  aboveskip=5pt, belowskip=5pt,
  escapeinside={<@}{@>},
  literate=
    {>=>}{{\hsop{>=>}}}3
    {>>=}{{\hsop{>\kern0pt>=}}}3
    {=<<}{{\hsop{=<\kern0pt<}}}3
    {<\$>}{{\hsop{<\$>}}}3
    {<*>}{{\hsop{<*>}}}3
    {<>}{{\hsop{<>}}}2
    {->}{{\hsop{$\rightarrow$}}}2
    {<-}{{\hsop{$\leftarrow$}}}2
    {=>}{{\hsop{$\Rightarrow$}}}2
    {::}{{\hsop{::}}}2
    {==}{{\hsop{==}}}2
    {0}{{\hsnum{0}}}1 {1}{{\hsnum{1}}}1 {2}{{\hsnum{2}}}1
    {3}{{\hsnum{3}}}1 {4}{{\hsnum{4}}}1 {5}{{\hsnum{5}}}1
    {6}{{\hsnum{6}}}1 {7}{{\hsnum{7}}}1 {8}{{\hsnum{8}}}1
    {9}{{\hsnum{9}}}1
    {ë}{{\"e}}1 {é}{{\'e}}1 {è}{{\`e}}1 {ï}{{\"i}}1,
}
\lstset{style=hs}
 
% inline code: same colors, plain text in dark blue
\lstdefinestyle{hsinline}{
  basicstyle=\ttfamily\color{kuldonker},
}
\newcommand{\hs}{\lstinline[style=hsinline]}
% voor code in argumenten van andere commando's (titels, tabellen, ...)
\newcommand{\code}[1]{{\ttfamily\color{kuldonker}#1}}
 
% antwoord dat pas op een latere overlay verschijnt (binnen listings)
\newcommand{\ans}[2]{\uncover<#1->{\color{codecomment}\itshape #2}}
 
% ---------------------------------------------------------------- tikz
\tikzset{
  obj/.style={font=\small},
  box/.style={draw=kuldonker, fill=kulblauw!10, rounded corners=2pt,
              inner sep=4pt, font=\small\ttfamily, text=kuldonker,
              minimum width=2.4cm},
  pow/.style={-{Stealth[length=5pt]}, thick, kuldonker},
}
\tikzcdset{arrow style=tikz, diagrams={>={Stealth[length=5pt]}}}
 
% ---------------------------------------------------------------- titel
\title{Sessie II: Haskell Type-classes}
\subtitle{PAL-sessie Declaratieve Talen}
\author{Yarne Ramakers, Yendric Van Roey}
\institute{KU Leuven}
\date{\today}
 
\begin{document}
 
\maketitle

\begin{frame}[standout]
  Zijn er vragen vanuit de oefenzittingen?
\end{frame}

\begin{frame}{Planning voor vandaag}
  \begin{enumerate}
    \item \textbf{Wat zijn typeclasses?}
    \item \textbf{Nuttige typeclasses}
    \begin{itemize}
      \item \textbf{Semigroepen en Monoiden}
      \item \textbf{Functoren}
      \item \textbf{Applicatives}
      \item \textbf{Foldable en Traversable}
      \item \textbf{Monad}
    \end{itemize}
    \item \textbf{Opdrachten}
  \end{enumerate}
\end{frame}

\section{Wat zijn typeclasses}

\begin{frame}[fragile]{Generics}
  Je wilt een functie \hs{show} die elk type omzet naar een \hs{String}:
\begin{lstlisting}
showInt  :: Int  -> String
showBool :: Bool -> String      -- en voor [Bool]? [[Int]]? ...
\end{lstlisting}
  \pause
  Kan het met een typevariabele? \hs{show :: a -> String}
  \pause
 
  Nee! Een functie die voor \alert{elke} \hs{a} werkt, weet niets over \hs{a}.
  Ze kan niet in de waarde kijken en moet altijd hetzelfde teruggeven.
  \pause
  \begin{itemize}
    \item \textbf{Parametrisch polymorfisme}: één implementatie voor alle types, bv.\ \hs{length}
    \item \textbf{Ad-hoc polymorfisme}: één naam, een \alert{eigen} implementatie per type\\
      $\Rightarrow$ \textbf{typeclasses}
  \end{itemize}
\end{frame}

\begin{frame}[fragile]{Wat is een typeclass?}
  \begin{block}{Definitie}
    Een \textbf{typeclass} is een verzameling functies (met bijhorende
    \alert{wetten}) die een type kan implementeren.
    Een type dat dat doet is een \textbf{instantie} van die klasse.
  \end{block}
  \pause
  Gelijkaardig aan interfaces uit OGP:
  \begin{columns}[T]
    \column{.5\textwidth}
\begin{lstlisting}[language=Java,keywords={public,interface,boolean},keywords={[2]Object},keywords={[3]equals}]
public interface Equals {
    public boolean equals(Object other);
}
\end{lstlisting}
    \column{.45\textwidth}
\begin{lstlisting}
class Eq a where
  (==) :: a -> a -> Bool
  (/=) :: a -> a -> Bool
\end{lstlisting}
  \end{columns}
  \pause
  \vspace{4pt}
  Maar sterker:
  \begin{itemize}
    \item je kan een instantie \alert{achteraf} toevoegen, los van de definitie van het type
    \item een klasse kan gaan over \alert{typeconstructors} zoals \hs{Maybe} of \hs{[]} (zie verder)
  \end{itemize}
\end{frame}

\begin{frame}[fragile]{Een instantie schrijven}
  \begin{columns}[T]
    \column{.48\textwidth}
\begin{lstlisting}
class Eq a where
  (==), (/=) :: a -> a -> Bool
  -- standaardimplementaties
  x /= y = not (x == y)
  x == y = not (x /= y)
  {-# MINIMAL (==) | (/=) #-}
\end{lstlisting}
    \column{.48\textwidth}
\begin{lstlisting}
data Color = Red | Green | Blue
 
instance Eq Color where
  Red   == Red   = True
  Green == Green = True
  Blue  == Blue  = True
  _     == _     = False
\end{lstlisting}
  \end{columns}
  \pause
  \vspace{6pt}
  \hs{MINIMAL}: je moet \hs{(==)} \alert{of} \hs{(/=)} implementeren, de andere krijg je gratis.
  \pause
 
  \vspace{6pt}
  Voor de standaardklassen kan de compiler het ook voor jou doen:
\begin{lstlisting}
data Color = Red | Green | Blue
  deriving (Show, Eq, Ord)
\end{lstlisting}
\end{frame}

\begin{frame}[fragile]{Constraints}
  Een functie kan eisen dat een type tot een klasse behoort:
\begin{lstlisting}
elem :: Eq a => a -> [a] -> Bool    -- voor elke a waarvoor Eq a bestaat
elem _ []       = False
elem x (y : ys) = x == y || elem x ys
\end{lstlisting}
  \pause
  Constraints kom je ook tegen bij klassen en instanties:
\begin{lstlisting}
class Eq a => Ord a where           -- elk Ord type is ook Eq (superklasse)
  compare :: a -> a -> Ordering
 
instance Eq a => Eq (Maybe a) where -- Maybe a is Eq als a Eq is
  Nothing == Nothing = True
  Just x  == Just y  = x == y
  _       == _       = False
\end{lstlisting}
\end{frame}
 
\begin{frame}[fragile]{Een typeclass is een contract}
  De compiler controleert de \alert{types}, maar niet de \alert{wetten}.
  Die zijn een belofte van wie de instantie schrijft.
  \pause
 
  \vspace{6pt}
  Voorbeeld: \hs{(==)} hoort een equivalentierelatie te zijn:
  \[
    \underbrace{\text{\code{x == x}}}_{\text{reflexief}}
    \qquad
    \underbrace{\text{\code{x == y}} \Rightarrow \text{\code{y == x}}}_{\text{symmetrisch}}
    \qquad
    \underbrace{\text{\code{x == y}} \land \text{\code{y == z}} \Rightarrow \text{\code{x == z}}}_{\text{transitief}}
  \]
  \pause
  Ook de standard library breekt dit soms:
\begin{lstlisting}
ghci> let nan = 0 / 0 :: Double
ghci> nan == nan
False                                -- niet reflexief!
\end{lstlisting}
  \pause
  Code die op de wetten rekent (bv.\ \hs{elem nan [nan]}) gedraagt zich dan raar.
\end{frame}

\begin{frame}[fragile]{Kinds: de types van types}
  Waarden hebben types, types hebben \textbf{kinds}:
\begin{lstlisting}
ghci> :kind Int
Int :: *                             -- een "gewoon" type
ghci> :kind Maybe
Maybe :: * -> *                      -- geef een type, krijg een type
ghci> :kind Either
Either :: * -> * -> *
\end{lstlisting}
  \hs{Maybe} op zich is geen type: er bestaan geen waarden van type \hs{Maybe}.
  \hs{Maybe Int} wel.
  \pause
 
  \vspace{6pt}
  Een typeclass kan over zulke \alert{typeconstructors} gaan:
\begin{lstlisting}
class Functor (f :: * -> *) where   -- f is Maybe, [], IO, ...
  fmap :: (a -> b) -> f a -> f b
\end{lstlisting}
\end{frame}

\begin{frame}[fragile]{Kleine quiz: kinds}
  Wat is de kind van:
\begin{lstlisting}
[]              -- <@\ans{2}{* -> *}@>
Maybe Int       -- <@\ans{2}{*}@>
Either String   -- <@\ans{2}{* -> *}@>
(,) Int         -- <@\ans{2}{* -> *}@>
(->) Int        -- <@\ans{2}{* -> *}@>
\end{lstlisting}
  \pause
  \pause
  Typeconstructors kan je dus ook \alert{partieel toepassen}, net als functies.
  \hs{Either String}, \hs{(,) Int} en \hs{(->) Int} kunnen daardoor
  instanties van \hs{Functor} zijn.
\end{frame}

\begin{frame}[fragile]{Documentatie vinden}
  \begin{itemize}
    \item \textbf{Hoogle}: zoek op naam \alert{of op type}:
      \href{https://hoogle.haskell.org}{\code{hoogle.haskell.org}}\\
      {\small \hs{(a -> b) -> [a] -> [b]} vindt \hs{map} en \hs{fmap}}
    \item \textbf{Hackage}: de documentatie van \code{base}, met per klasse
      de \alert{wetten} en de minimale definitie:
      \href{https://hackage.haskell.org/package/base}{\code{hackage.haskell.org/package/base}}
    \item \textbf{GHCi}: \hs{:type}, \hs{:kind} en \hs{:info}
  \end{itemize}
\begin{lstlisting}
ghci> :info Functor
type Functor :: (* -> *) -> Constraint
class Functor f where
  fmap :: (a -> b) -> f a -> f b
  (<$) :: a -> f b -> f a
  {-# MINIMAL fmap #-}
instance Functor Maybe
instance Functor []
...
\end{lstlisting}
\end{frame}



\section{Semigroepen en monoïden}

\begin{frame}{Combineren}
  We combineren constant waarden: $3 + 4$, \code{"ab" ++ "cd"},
  \code{max 3 7}, foutmeldingen samenvoegen, \dots
  \pause
  \begin{block}{Definitie}
    Een \textbf{semigroep} $(S, \cdot)$ is een verzameling $S$ met een binaire
    operatie $\cdot : S \times S \to S$ die \alert{associatief} is:
    \[ \forall x, y, z \in S : (x \cdot y) \cdot z = x \cdot (y \cdot z) \]
    Een \textbf{monoïde} $(M, \cdot, e)$ is een semigroep met een
    \alert{neutraal element} $e \in M$:
    \[ \forall x \in M : e \cdot x = x = x \cdot e \]
  \end{block}
  \pause
  Voorbeelden: $(\mathbb{N}, +, 0)$, \ $(\mathbb{Z}, \times, 1)$, \
  (strings, concatenatie, $\varepsilon$).\\
  Geen voorbeeld: $(\mathbb{Z}, -)$, want $(5 - 3) - 1 = 1 \neq 3 = 5 - (3 - 1)$.
\end{frame}
 
\begin{frame}[fragile]{Semigroup en Monoid in Haskell}
\begin{lstlisting}
class Semigroup a where
  (<>) :: a -> a -> a
  -- wet: (x <> y) <> z == x <> (y <> z)
 
class Semigroup a => Monoid a where
  mempty :: a
  -- wetten: mempty <> x == x
  --         x <> mempty == x
\end{lstlisting}
  \pause
  \begin{columns}[T]
    \column{.4\textwidth}
\begin{lstlisting}
instance Semigroup [a] where
  (<>) = (++)
 
instance Monoid [a] where
  mempty = []
\end{lstlisting}
    \column{.55\textwidth}
    Gratis erbij:
\begin{lstlisting}
mconcat :: Monoid a => [a] -> a
mconcat ["ab", "c", "de"]  == "abcde"
mconcat []                 == ""
\end{lstlisting}
  \end{columns}
\end{frame}
 
\begin{frame}{Voorbeelden}
  \centering
  \begin{tabular}{lllll}
    \toprule
    \textbf{Type} & \textbf{\code{<>}} & \textbf{\code{mempty}} & \textbf{Klasse} \\
    \midrule
    \code{[a]}         & \code{++}            & \code{[]}    & Monoid \\
    \code{Sum Int}     & \code{+}             & \code{0}     & Monoid \\
    \code{Product Int} & \code{*}             & \code{1}     & Monoid \\
    \code{All}         & \code{\&\&}          & \code{True}  & Monoid \\
    \code{Any}         & \code{||}            & \code{False} & Monoid \\
    \code{Max a}       & \code{max}           & \code{minBound}, enkel als \code{Bounded a} & Semigroup (Monoid) \\
    \code{NonEmpty a}  & \code{++}            & --- (een lege lijst is niet toegelaten!) & Semigroup \\
    \code{Int} met \code{-} & --- & --- & geen van beide \\
    \bottomrule
  \end{tabular}
 
  \vspace{10pt}
  Bestaat er wel een \alert{zinnig} neutraal element?
\end{frame}
 
\begin{frame}{Kleine quiz: semigroep, monoïde of geen van beide?}
  \begin{enumerate}
    \item \code{Double} met het gemiddelde \code{\textbackslash x y -> (x + y) / 2}
      \only<2->{\\ \alert{geen van beide}: $\text{gem}(\text{gem}(0, 0), 4) = 2$, maar $\text{gem}(0, \text{gem}(0, 4)) = 1$}
    \item \code{Bool} met \code{xor}
      \only<3->{\\ \alert{monoïde}: associatief, met neutraal element \code{False}}
    \item \code{Maybe a} met ``neem de eerste \code{Just}''
      \only<4->{\\ \alert{monoïde}, met neutraal element \code{Nothing} (dit is \code{First a})}
    \item functies \code{a -> a} met compositie \code{(.)}
      \only<5->{\\ \alert{monoïde}: \code{(.)} is associatief en \code{id} is neutraal.}
  \end{enumerate}
\end{frame}

\begin{frame}[fragile]{Eén type, meerdere monoïden}
  \hs{Int} heeft twee voor de hand liggende monoïden: $(+, 0)$ en $(\times, 1)$.
  Welke instantie kies je?
\begin{lstlisting}
newtype Sum a = Sum { getSum :: a }          -- uit Data.Monoid
 
instance Num a => Semigroup (Sum a) where
  Sum x <> Sum y = Sum (x + y)
instance Num a => Monoid (Sum a) where
  mempty = Sum 0
\end{lstlisting}
  \pause
\begin{lstlisting}
getSum     (mconcat (map Sum     [1 .. 4]))   == 10
getProduct (mconcat (map Product [1 .. 4]))   == 24
\end{lstlisting}
\end{frame}

\begin{frame}[fragile]{Bestaat er wel een neutraal element?}
  Niet elke operatie heeft een neutraal element. \hs{max} heeft er geen voor elk type:
\begin{lstlisting}
ghci> maximum []
*** Exception: Prelude.maximum: empty list
\end{lstlisting}
  \pause
  Combineer je enkel als er \alert{minstens één} element is? Dan volstaat een semigroep:
\begin{lstlisting}
sconcat :: Semigroup a => NonEmpty a -> a
\end{lstlisting}
  \pause
  \begin{alertblock}{Vuistregel}
    Begin met \code{Semigroup}. Voeg \code{mempty} enkel toe als er in je domein
    een betekenisvolle ``doe niets''-waarde bestaat.
  \end{alertblock}
\end{frame}

\section{Functoren}

\begin{frame}[fragile]{Functor: de wiskunde}
  \begin{columns}[c]
    \column{.6\textwidth}
    \begin{block}{Definitie}
      Een \textbf{functor} $F : \mathcal{C} \to \mathcal{D}$ beeldt
      \begin{itemize}
        \item elk object $A$ af op een object $F\,A$
        \item elk morfisme $f : A \to B$ af op $F\,f : F\,A \to F\,B$
      \end{itemize}
      zodat de structuur bewaard blijft:
      \[
        F\,\mathrm{id}_A = \mathrm{id}_{F A}
        \qquad
        F\,(g \circ f) = F\,g \circ F\,f
      \]
    \end{block}
    Een \textbf{endofunctor} gaat van een categorie naar zichzelf:
    \alert{Hask $\to$ Hask}.
    \column{.38\textwidth}
    \centering
    \begin{tikzcd}[row sep=huge, column sep=large]
      A \arrow[r, "f"] \arrow[d, mapsto, dashed, "F"']
        & B \arrow[d, mapsto, dashed, "F"] \\
      F\,A \arrow[r, "F\,f"']
        & F\,B
    \end{tikzcd}
  \end{columns}
\end{frame}

\begin{frame}[fragile]{Functor in Haskell}
  Een endofunctor op Hask bestaat uit twee delen:
  \begin{itemize}
    \item op objecten: een \alert{typeconstructor} \hs{f :: * -> *} \hfill \hs{Int} $\mapsto$ \hs{Maybe Int}
    \item op morfismen: \alert{\hs{fmap}} \hfill \hs{a -> b} $\mapsto$ \hs{f a -> f b}
  \end{itemize}
\begin{lstlisting}
class Functor f where
  fmap :: (a -> b) -> (f a -> f b)       -- "til een functie op"
-- wetten (= de definitie van een functor!)
fmap id      == id
fmap (g . h) == fmap g . fmap h
\end{lstlisting}
  \pause
  \hs{fmap} verandert de \alert{waarden}, nooit de \alert{structuur}.
  Twee keer mappen = één keer mappen met de samenstelling.\\[4pt]
  \hs{(<$>)} is \hs{fmap} als infix-operator: \hs{f <$> x == fmap f x}
\end{frame}
 
\begin{frame}[fragile]{Functor-instanties}
  \begin{columns}[T]
    \column{.5\textwidth}
\begin{lstlisting}
instance Functor Maybe where
  fmap _ Nothing  = Nothing
  fmap f (Just x) = Just (f x)
 
instance Functor [] where
  fmap = map
\end{lstlisting}
    \column{.48\textwidth}
\begin{lstlisting}
data Tree a = Leaf a
            | Node (Tree a) (Tree a)
 
instance Functor Tree where
  fmap f (Leaf x)   = Leaf (f x)
  fmap f (Node l r) =
    Node (fmap f l) (fmap f r)
\end{lstlisting}
  \end{columns}
  \pause
\begin{lstlisting}
fmap (+ 1) (Just 5)                == Just 6
fmap (+ 1) Nothing                 == Nothing
length <$> ["Hello", "World", "!"] == [5, 5, 1]
show <$> Node (Leaf 1) (Leaf 2)    == Node (Leaf "1") (Leaf "2")
\end{lstlisting}
\end{frame}

\begin{frame}[fragile]{Kleine quiz: functoren}
  Wat is het resultaat?
\begin{lstlisting}
fmap (+ 1) (Right 3)       -- <@\ans{2}{Right 4}@>
fmap (+ 1) (Left "fout")   -- <@\ans{2}{Left "fout"}@>
fmap (+ 1) (3, 4)          -- <@\ans{3}{(3, 5)}@>
fmap (* 2) (+ 1) 3         -- <@\ans{4}{8}@>
\end{lstlisting}
  \only<2>{\hs{Either e} is een functor in het \alert{laatste} argument}
  \only<3>{\hs{(,) a} heeft kind \hs{* -> *}: het eerste element hoort bij de \alert{structuur}, enkel het tweede wordt gemapt.}
  \only<4>{\hs{(->) r} is ook een functor: \hs{fmap} voor functies is \alert{compositie}.
    \hs{fmap (* 2) (+ 1) == (* 2) . (+ 1)}}
\end{frame}
 
\begin{frame}[fragile]{Kleine quiz: is dit een geldige functor?}
  \begin{columns}[T]
    \column{.48\textwidth}
\begin{lstlisting}
instance Functor [] where
  fmap f xs = reverse (map f xs)
\end{lstlisting}
    \column{.48\textwidth}
\begin{lstlisting}
data Counted a = Counted Int a
 
instance Functor Counted where
  fmap f (Counted n x) =
    Counted (n + 1) (f x)
\end{lstlisting}
  \end{columns}
  \pause
  Beide \alert{typechecken}, maar geen van beide is een functor:
\begin{lstlisting}
fmap id [1, 2, 3]        == [3, 2, 1]        -- /= id [1, 2, 3]
fmap id (Counted 0 'x')  == Counted 1 'x'    -- /= id (Counted 0 'x')
\end{lstlisting}
  \pause
  Beide veranderen de \alert{structuur} en dus schenden de wetten van een functor!
\end{frame}

\section{Applicatives}

\begin{frame}[fragile]{De grens van Functor}
  Wat als je een functie met \alert{twee} argumenten hebt?
\begin{lstlisting}
ghci> :type fmap (+) (Just 3)
fmap (+) (Just 3) :: Maybe (Int -> Int)
\end{lstlisting}
  Nu zit er een \alert{functie} in een context, en die wil je toepassen op \hs{Just 4}
  \pause
\begin{lstlisting}
class Functor f => Applicative f where
  pure  :: a -> f a                     -- waarde in minimale context
  (<*>) :: f (a -> b) -> f a -> f b     -- ingepakte functie toepassen
\end{lstlisting}
  \pause
\begin{lstlisting}
(+)                        :: Int -> Int -> Int
(+) <$> Just 3             :: Maybe (Int -> Int)
(+) <$> Just 3 <*> Just 4  :: Maybe Int           -- Just 7
\end{lstlisting}
\end{frame}

\begin{frame}[fragile]{Applicative-instanties}
  \begin{columns}[T]
    \column{.48\textwidth}
\begin{lstlisting}
instance Applicative Maybe where
  pure = Just
  Just f <*> Just x = Just (f x)
  _      <*> _      = Nothing
\end{lstlisting}
    \column{.48\textwidth}
\begin{lstlisting}
instance Applicative [] where
  pure x    = [x]
  fs <*> xs = [f x | f <- fs, x <- xs]
\end{lstlisting}
  \end{columns}
  \pause
\begin{lstlisting}
(+) <$> Just 3 <*> Just 4         == Just 7
(+) <$> Just 3 <*> Nothing        == Nothing
(+) <$> [1, 2] <*> [10, 20]       == [11, 21, 12, 22]
(,) <$> [1, 2] <*> "ab"           == [(1,'a'), (1,'b'), (2,'a'), (2,'b')]
\end{lstlisting}
  \pause
  \begin{itemize}
    \item \hs{Maybe}: één \hs{Nothing} en het geheel faalt
    \item lijsten = \alert{cartesisch product}
  \end{itemize}
\end{frame}

\begin{frame}[fragile]{Applicative wetten}
\begin{lstlisting}
pure id <*> v              == v                  -- identiteit
pure f <*> pure x          == pure (f x)         -- homomorfisme
u <*> pure y               == prue ($ y) <*> u   -- interchange
pure (.) <*> u <*> v <*> w == u <*> (v <*> w)    -- compositie
\end{lstlisting}
en consistent met \hs|Functor| \hs|fmap f x == pure f <*> x|
\pause 
\begin {itemize}
\item \hs|pure| voegt \alert{geen effecten} toe (neutrale element)
\item effecten combineren \alert{associatief}
\end {itemize}
Klinkt bekend?
\pause 
Een Applicative is een monoïde over contexten (effecten, ...)
\end{frame}

\begin{frame}[fragile]{Handige Applicative-functies}
\begin{lstlisting}[morekeywords={[3]{replicateM,when}}]
liftA2     :: Applicative f => (a -> b -> c) -> f a -> f b -> f c
replicateM :: Applicative f => Int -> f a -> f [a]      -- Control.Monad
when       :: Applicative f => Bool -> f () -> f ()     -- Control.Monad
liftA2 g x y = g <$> x <*> y
when b act   = if b then act else pure ()
\end{lstlisting}
  \pause
\begin{lstlisting}[morekeywords={[3]{replicateM,when}}]
liftA2 (+) (Just 3) (Just 4)     == Just 7
liftA2 (,) [1, 2] "ab"           == [(1,'a'), (1,'b'), (2,'a'), (2,'b')]
replicateM 2 "ab"                == ["aa", "ab", "ba", "bb"]
replicateM 3 getLine             -- lees 3 regels: IO [String]
when verbose (putStrLn "klaar")  -- niet verbose? dan pure ()
\end{lstlisting}
\end{frame}

\begin{frame}[fragile]{Applicatives zijn statisch}
\begin{lstlisting}
User <$> validateName n <*> validateAge a
\end{lstlisting}
\begin{itemize}
\item de effecten zijn \alert{onafhankelkijk}, \hs|validateAge| weet niets over het resultaat van \hs|validateName|
\item de vorm van de berkeneing ligt vast \alert{voor} er iets uitegevoerd wordt
\end{itemize}
\pause 
Dus je kan hier mee alle effecten uitvoeren (Monad kan eerder stoppen) en deze mogen in \alert{parallel} lopen.
\pause 
Maar je kan \alert{niet} zeggen "als eerste resultaat > 0", doen dan dit.
\end{frame}

\section{Foldable}

\begin{frame}[fragile]{Foldable: een structuur samenvatten}
  \hs{Foldable} vat alle waarden van een structuur samen tot één resultaat
  \pause
\begin{lstlisting}
class Foldable t where
  foldMap :: Monoid m => (a -> m) -> t a -> m
  foldr   :: (a -> b -> b) -> b -> t a -> b
  {-# MINIMAL foldMap | foldr #-}
\end{lstlisting}
  \pause
  \hs{foldMap f}: zet elk element om naar een \alert{monoïde} en combineer deze met \hs{<>}.
  \begin{itemize}
    \item \hs{<>} combineert de elementen, \hs{mempty} is de fallback voor een lege structuur
    \item hoe je groepeert maakt niet uit. Het resultaat is altijd hetzelfde want \hs{<>} is associatief
  \end{itemize}
\end{frame}

\begin{frame}[fragile]{\texttt{foldMap} in actie}
\begin{lstlisting}
foldMap Sum            [1, 2, 3, 4]   == Sum 10
foldMap Product        [1, 2, 3, 4]   == Product 24
foldMap (Any . even)   [1, 3, 4]      == Any True
foldMap (\x -> [x, x]) [1, 2]         == [1, 1, 2, 2]
foldMap Sum            []             == Sum 0          -- mempty
\end{lstlisting}
  \pause
  De monoïde bepaalt wat er berekend wordt:
\begin{lstlisting}
sum    = getSum . foldMap Sum
length = getSum . foldMap (const (Sum 1))
elem x = getAny . foldMap (\y -> Any (x == y))
toList = foldMap (\x -> [x])
\end{lstlisting}
\end{frame}

\begin{frame}[fragile]{Een Foldable-instantie schrijven}
  \begin{columns}[T]
    \column{.5\textwidth}
\begin{lstlisting}
data Tree a = Leaf a
            | Node (Tree a) (Tree a)

instance Foldable Tree where
  foldMap f (Leaf x)   = f x
  foldMap f (Node l r) =
    foldMap f l <> foldMap f r
\end{lstlisting}
    \column{.46\textwidth}
    Je kiest \alert{waar} de elementen zitten en in welke structuur,
    maar hoe ze gecombineerd worden kiest de gebruiker via zijn \alert{monoïde}
  \end{columns}
  \pause
\begin{lstlisting}
t = Node (Leaf 1) (Node (Leaf 2) (Leaf 3))
sum t  == 6        length t  == 3        maximum t == 3
elem 2 t == True   toList t  == [1, 2, 3]
\end{lstlisting}
  Gratis erbij: \hs{sum}, \hs{product}, \hs{length}, \hs{elem}, \hs{maximum},
  \hs{null}, \hs{toList}, \hs{foldr}, \dots
\end{frame}

\begin{frame}[fragile]{Kleine quiz: Foldable}
  Wat is het resultaat?
\begin{lstlisting}
length (Just 'x')      -- <@\ans{2}{1}@>
sum Nothing            -- <@\ans{2}{0}@>
length (3, 4)          -- <@\ans{3}{1}@>
elem 3 (3, 4)          -- <@\ans{3}{False}@>
null (Left "fout")     -- <@\ans{3}{True}@>
\end{lstlisting}
  \only<2>{\hs{Maybe} is een container met \alert{0 of 1} elementen.}
  \only<3>{cf. \hs{Functor}: \hs{(,) a} en \hs{Either e} hebben kind \hs{* -> *}.
    \\Het laatste typeargument is de inhoud.}
\end{frame}

\section{Traversable}

\begin{frame}[fragile]{Het probleem}
  Een functie die kan \alert{falen} (of ander effect), toegepast op elk element:
\begin{lstlisting}
import Text.Read (readMaybe)    -- readMaybe :: Read a => String -> Maybe a

fmap readMaybe ["1", "2", "3"]  == [Just 1, Just 2, Just 3]  -- [Maybe Int]
\end{lstlisting}
  \pause
  We willen ``\alert{alles} is gelukt, of het geheel faalt'': \hs{Maybe [Int]}.\\
  We willen dus de effecten van in de structuur naar buiten krijgen.
  \pause
\begin{lstlisting}
fmap     :: (a -> f b) -> t a -> t (f b)    -- effecten zitten binnenin
traverse :: (a -> f b) -> t a -> f (t b)    -- effecten naar buiten

traverse readMaybe ["1", "2", "3"]  == Just [1, 2, 3]
traverse readMaybe ["1", "x", "3"]  == Nothing
\end{lstlisting}
\end{frame}

\begin{frame}[fragile]{Traversable}
\begin{lstlisting}
class (Functor t, Foldable t) => Traversable t where
  traverse  :: Applicative f => (a -> f b) -> t a -> f (t b)
  sequenceA :: Applicative f => t (f a) -> f (t a)
  {-# MINIMAL traverse | sequenceA #-}

traverse f = sequenceA . fmap f
sequenceA  = traverse id
\end{lstlisting}
  \pause
  Idee: \hs{fmap} met effecten. Doorloop de structuur, doe bij elk element een effect, behoud de structuur en combineer alle effecten tot één.
  \pause

  \vspace{4pt}
  \hs{sequenceA} brengt uw effect naar buiten en behoudt de structuur:
\begin{lstlisting}
sequenceA [Just 1, Just 2]    == Just [1, 2]
sequenceA [Just 1, Nothing]   == Nothing
\end{lstlisting}
\end{frame}

\begin{frame}[fragile]{Een Traversable-instantie schrijven}
  \begin{columns}[T]
    \column{.44\textwidth}
\begin{lstlisting}
instance Traversable [] where
  traverse f []       = pure []
  traverse f (x : xs) =
    (:) <$> f x <*> traverse f xs
\end{lstlisting}
    \column{.54\textwidth}
\begin{lstlisting}
instance Traversable Tree where
  traverse f (Leaf x)   = Leaf <$> f x
  traverse f (Node l r) =
    Node <$> traverse f l <*> traverse f r
\end{lstlisting}
  \end{columns}
  \pause
  Elke constructor \hs{C} wordt \hs{C <$> ... <*> ...}
  \begin{itemize}
    \item de constructor bouwt de structuur opnieuw op
    \item \hs{<*>} combineert de effecten
  \end{itemize}
\end{frame}

\begin{frame}[fragile]{Stap voor stap}
\begin{lstlisting}
safeHalve :: Int -> Maybe Int
safeHalve x = if even x then Just (x `div` 2) else Nothing
\end{lstlisting}
  \pause
  Uitgeschreven:
\begin{lstlisting}
traverse safeHalve [2, 4, 6]
  = (:) <$> Just 1 <*> ((:) <$> Just 2 <*> ((:) <$> Just 3 <*> pure []))
\end{lstlisting}
  \pause
  Combineer van binnen naar buiten (\hs{(:) <$> Just x == Just (x :)}):
\begin{lstlisting}
Just (3 :) <*> Just []         == Just [3]
Just (2 :) <*> Just [3]        == Just [2, 3]
Just (1 :) <*> Just [2, 3]     == Just [1, 2, 3]
\end{lstlisting}
\end{frame}

\begin{frame}[fragile]{Wat als één stap faalt?}
\begin{lstlisting}
traverse safeHalve [2, 3, 6]
  = (:) <$> Just 1 <*> ((:) <$> Nothing <*> ((:) <$> Just 3 <*> pure []))
\end{lstlisting}
  \pause
\begin{lstlisting}
Just (3 :)  <*> Just []        == Just [3]
Nothing     <*> Just [3]       == Nothing
Just (1 :)  <*> Nothing        == Nothing
\end{lstlisting}
  \pause
  Eén \hs{Nothing} verspreidt zich naar buiten via elke \hs{<*>}.\\
  Je krijgt dit volledig gratis uit \hs{<*>} van (in dit geval) \hs{Maybe}.
\end{frame}

\begin{frame}[fragile]{Zelfde \texttt{traverse}, andere \texttt{<*>}}
\begin{lstlisting}
-- Either e: de eerste fout wint
check x = if x > 0 then Right x else Left ("fout: " ++ show x)
traverse check [1, 2, 3]         == Right [1, 2, 3]
traverse check [1, -2, -3]       == Left "fout: -2"

-- lijsten: alle combinaties
traverse (\x -> [x, -x]) [1, 2]  == [[1, 2], [1, -2], [-1, 2], [-1, -2]]

-- IO: voer de acties in volgorde uit
traverse print [1, 2, 3]         -- print 1, 2, 3 en geeft [(), (), ()]
\end{lstlisting}

  \pause

  {\small Reulstaat niet nodig? \hs{traverse_} / \hs{mapM_}
  vragen enkel \hs{Foldable} (resultaat wordt \hs{()}))}
\end{frame}

\begin{frame}[fragile]{Kleine quiz: Traversable}
  Wat is het resultaat?
\begin{lstlisting}
sequenceA [Just 1, Just 2]               -- <@\ans{2}{Just [1, 2]}@>
sequenceA [Just 1, Nothing]              -- <@\ans{2}{Nothing}@>
sequenceA [[1, 2], [3, 4]]               -- <@\ans{3}{[[1, 3], [1, 4], [2, 3], [2, 4]]}@>
sequenceA ([] :: [Maybe Int])               -- <@\ans{4}{Just []}@>
traverse readMaybe ("leeftijd", "21")       -- <@\ans{5}{Just ("leeftijd", 21)}@>
traverse safeHalve (Node (Leaf 2) (Leaf 5)) -- <@\ans{5}{Nothing}@>
\end{lstlisting}
  \only<3>{Lijst-Applicative = \alert{cartesisch product}: elke keuze met elke keuze.}
  \only<4>{Lege structuur: geen effecten om te combineren, enkel \hs{pure []}.}
\end{frame}

\begin{frame}[fragile]{Traversable is sterker dan Functor en Foldable}
  Met de juiste Applicative krijg je \hs{fmap} en \hs{foldMap} terug:
\begin{lstlisting}
newtype Identity a = Identity { runIdentity :: a }  -- effect dat niets doet
newtype Const m a  = Const { getConst :: m }        -- verzamelt enkel een m

instance Monoid m => Applicative (Const m) where
  pure _              = Const mempty
  Const x <*> Const y = Const (x <> y)
\end{lstlisting}
  \pause
\begin{lstlisting}
fmap    f = runIdentity . traverse (Identity . f)   -- geen effect
foldMap f = getConst    . traverse (Const . f)      -- enkel effect
\end{lstlisting}
\end{frame}


\begin{frame}{Samenvatting}
  \centering
  \begin{tabular}{lllll}
    \toprule
    \textbf{Klasse} & \textbf{Functie} & \textbf{Waarden} & \textbf{Vorm} & \textbf{Gebruikt} \\
    \midrule
    \code{Functor}     & \code{fmap}     & veranderd                 & behouden & --- \\
    \code{Foldable}    & \code{foldMap}  & samengevat                & weg      & \code{Monoid} \\
    \code{Traversable} & \code{traverse} & veranderd, \alert{met effect} & behouden & \code{Applicative} \\
    \bottomrule
  \end{tabular}

  \vspace{12pt}
  \begin{itemize}
    \item \textbf{Functor}: verander elke waarde, behoud de vorm
    \item \textbf{Foldable}: plet alle waarden tot één, verlies de vorm
    \item \textbf{Traversable}: verander elke waarde \alert{met een effect}, behoud de vorm
      en combineer de effecten tot één
  \end{itemize}
\end{frame}

\section{Monad}

\begin{frame}[fragile]{Het probleem}
\begin{lstlisting}
safeHalve :: Int -> Maybe Int
safeHalve x
  | even x    = Just (x `div` 2)
  | otherwise = Nothing
\end{lstlisting}

Hoe twee keer halveren?
\pause
\begin{lstlisting}
fmap safeHalve (Just 8)                    :: Maybe (Maybe Int)
fmap (fmap safeHalve) (fmap safeHalve ...) :: Maybe (Maybe (Maybe Int))
\end{lstlisting}
\pause 
\begin{itemize}
\item \hs|fmap| met een functie \hs|a -> m b| geeft \alert{geneste} contexten
\item \hs|<*>| helpt niet, of de tweede stap faalt hangt af van het \alert{resultaat} van de eerste
\end{itemize}
We willen een functie \hs|join :: Maybe (Maybe a) -> Maybe a|
\end{frame}

\begin{frame}[fragile]{Monad}
\begin{lstlisting}
class Applicative m => Monad m where
  (>>=)  :: m a -> (a -> m b) -> m b      -- "bind"
  return :: a -> m a                      -- = pure

join :: Monad m => m (m a) -> m a
join mma = mma >>= id
m >>= f = join (fmap f m)
\end{lstlisting}
\pause 
\begin{lstlisting}
Just 8 >>= safeHalve >>= safeHalve    == Just 2
Just 6 >>= safeHalve >>= safeHalve    == Nothing
\end{lstlisting}
\end{frame}

\begin{frame}[fragile]{Naast elkaar}
\begin{lstlisting}
(<$>) :: Functor f       =>   (a ->   b) -> f a -> f b
(<*>) :: Applicative f   => f (a ->   b) -> f a -> f b
(=<<) :: Monad f         =>   (a -> f b) -> f a -> f b
\end{lstlisting}
(\hs|=<<| is \hs|>>=| met de argumentne omgedraaid)
\pause
\begin{itemize}
\item \textbf{Functor}: een gewone functie, de context blijft wat ze is
\item \textbf{Applicative}: functie zit in de context, contexten worden \alert{gecombineerd}
\item \textbf{Monad}: de functie \alert{maakt} een nieuwe context, obv de waarde
\end{itemize}
Elke monad is een applicative, elke applicative is een functor, maar niet omgekeerd
\end{frame}

\begin{frame}[fragile]{Monad-instanties}
  \begin{columns}[T]
    \column{.48\textwidth}
\begin{lstlisting}
instance Monad Maybe where
  Nothing >>= _ = Nothing
  Just x  >>= f = f x
\end{lstlisting}
\begin{lstlisting}
join (Just (Just 4)) == Just 4
join (Just Nothing)  == Nothing
\end{lstlisting}
    \column{.48\textwidth}
\begin{lstlisting}
instance Monad [] where
  xs >>= f = concatMap f xs
\end{lstlisting}
\begin{lstlisting}
join [[1, 2], [], [3]] == [1, 2, 3]
\end{lstlisting}
  \end{columns}
  \pause
\begin{lstlisting}
[1, 2, 3] >>= \x -> [x, -x]   == [1, -1, 2, -2, 3, -3]
\end{lstlisting}
  \pause
  \begin{itemize}
    \item \hs{Maybe}stop bij de eerste \hs{Nothing} (\alert{early exit})
    \item lijsten proberen alle mogelijkheden (\alert{non-determinisme})
  \end{itemize}
\end{frame}

\begin{frame}[fragile]{do-notatie}
  \hs{do} is syntactic sugar voor \hs{>>=}:
  \begin{columns}[T]
    \column{.48\textwidth}
\begin{lstlisting}
do x <- m
   y <- k x
   return (x, y)
\end{lstlisting}
    \column{.48\textwidth}
\begin{lstlisting}
m   >>= \x ->
k x >>= \y ->
return (x, y)
\end{lstlisting}
  \end{columns}
  \pause
\begin{lstlisting}
data User    = User    { address :: Maybe Address }
data Address = Address { city    :: Maybe String }
 
getCity :: User -> Maybe String
getCity user = do addr <- address user     -- Nothing? dan stopt het hier
                  city addr
\end{lstlisting}
  \pause
  Let op: \hs{return} stopt de functie \alert{niet}! Het steekt enkel een waarde in de context.
\end{frame}

\begin{frame}[fragile]{Monad wetten}
\begin{lstlisting}
return a >>= f     == f a                     -- linker-identiteit
m >>= return       == m                       -- rechter-idnetiteit
(m >>= f) >>= g    == m >>= (\x -> f x >>= g) -- associativiteit
\end{lstlisting}
\pause 
\alert{Kleisli-compositie} stelt functies \hs|a -> m b| samen
\begin{lstlisting}
import Control.Monad ((>=>))
(>=>) :: Monad m => (a -> m b) -> (b -> m c) -> (a -> m c)
(f >=> g) x = f x >>= g
\end{lstlisting}
\begin{lstlisting}
return >=> f       == f
f >=> return       == f
(f >=> g) >=> h    == f >=> (g >=> h)
\end{lstlisting}
associatief met een neutraal element (\alert{monoidewetten})
\end{frame}

\begin{frame}[plain]
  \makebox[\linewidth]{%
    \includegraphics[width=\paperwidth,height=\paperheight,keepaspectratio]{cz76ci98ct0d1.png}%
  }
\end{frame}

\begin{frame}{Een monade is een monoide over de categorie van de C-endofunctoren}
\begin{itemize}
    \item \textbf{De categorie}: objecten zijn \alert{endofunctoren} (\code{Maybe}, \code{[]}, \code{IO}, \dots),
      morfismen zijn \alert{natuurlijke transformaties} \code{forall a. f a -> g a}\\
      {\small bv.\ \code{maybeToList :: Maybe a -> [a]}: werkt voor elke \code{a}, kijkt niet naar de inhoud}
    \pause
    \item \textbf{``Vermenigvuldigen''} van objecten: functoren \alert{samenstellen}
      ($m \circ m$ geeft \code{m (m a)}), met als eenheid de identiteitsfunctor $\mathrm{Id}$
    \pause
    \item \textbf{Een monoïde} in die categorie: een object $m$ met
  \end{itemize}
  \pause
  \centering
  \begin{tabular}{lll}
    \toprule
    & \textbf{Monoïde} $(\mathbb{Z}, +, 0)$ & \textbf{Monad} $m$ \\
    \midrule
    element      & een getal & een functor \code{m} \\
    combineren   & $+ : \mathbb{Z} \times \mathbb{Z} \to \mathbb{Z}$ & $\mu$ = \code{join :: m (m a) -> m a} \\
    neutraal     & $0$ & $\eta$ = \code{return :: a -> m a} \\
    \bottomrule
  \end{tabular}
\end{frame}

\begin{frame}[fragile]{De monade wetten zijn monoidewetten}
\begin{columns}[c]
    \column{.5\textwidth}
    \centering
    \textbf{Associativiteit}\\[4pt]
    \begin{tikzcd}[column sep=large, row sep=large]
      \texttt{m (m (m a))} \arrow[r, "\texttt{fmap join}"] \arrow[d, "\texttt{join}"']
        & \texttt{m (m a)} \arrow[d, "\texttt{join}"] \\
      \texttt{m (m a)} \arrow[r, "\texttt{join}"']
        & \texttt{m a}
    \end{tikzcd}\\[4pt]
    {\small $(a + b) + c = a + (b + c)$}
    \column{.5\textwidth}
    \centering
    \textbf{Identiteit}\\[4pt]
    \begin{tikzcd}[column sep=large, row sep=large]
      \texttt{m a} \arrow[r, "\texttt{return}"] \arrow[dr, equal]
        & \texttt{m (m a)} \arrow[d, "\texttt{join}"]
        & \texttt{m a} \arrow[l, "\texttt{fmap return}"'] \arrow[dl, equal] \\
      & \texttt{m a} &
    \end{tikzcd}\\[4pt]
    {\small $0 + a = a = a + 0$}
  \end{columns}
  Een monade is een manier om effecten monoïdisch samen te stellen.
\end{frame}

\begin{frame}{Kleine quiz: wat heb je minimaal nodig?}
  Functor, Applicative of Monad?
  \begin{enumerate}
    \item alle prijzen in een \code{Maybe [Double]} met 21\% btw verhogen
      \only<2->{\hfill \alert{Functor}}
    \item een \code{User} bouwen uit drie onafhankelijk gevalideerde velden
      \only<3->{\hfill \alert{Applicative}}
    \item een lijst van URL's downloaden tot \code{IO [String]}
      \only<4->{\hfill \alert{Applicative} (\code{traverse})}
    \item een bestandsnaam inlezen en dan dat bestand openen
      \only<5->{\hfill \alert{Monad}}
    \item een parser die eerst een getal $n$ leest en dan $n$ tekens
      \only<6->{\hfill \alert{Monad}}
  \end{enumerate}
\end{frame}

\begin{frame}{Je gebruikt al monads}
  \centering
  \begin{tabular}{llll}
    \toprule
    \textbf{Taal} & \textbf{Context} & \textbf{\code{return}} & \textbf{\code{>>=}} \\
    \midrule
    JavaScript & \code{Promise}           & \code{Promise.resolve} & \code{.then} \\
    Rust       & \code{Option}, \code{Result} & \code{Some}, \code{Ok} & \code{and\_then}, \code{?} \\
    Java       & \code{Optional}          & \code{Optional.of}     & \code{flatMap} \\
    C\#        & \code{IEnumerable}       & ---                    & \code{SelectMany} (LINQ) \\
    Python     & lijsten                  & \code{[x]}             & list comprehensions \\
    \bottomrule
  \end{tabular}
 
  \vspace{10pt}
\end{frame}

\begin{frame}[fragile]{Toestand zonder mutatie}
  Zuivere functies kunnen niets aanpassen. Toestand moet je dus \alert{doorgeven}:
\begin{lstlisting}
-- nummer elk blad van een boom, van links naar rechts
label :: Tree a -> Int -> (Tree (Int, a), Int)
label (Leaf x)   n = (Leaf (n, x), n + 1)
label (Node l r) n =
  let (l', n1) = label l n
      (r', n2) = label r n1       -- n1, niet n! makkelijk om fout te doen
  in  (Node l' r', n2)
\end{lstlisting}
  \pause
  Kijk naar het type: \hs{Int -> (Tree (Int, a), Int)} heeft de vorm \hs{s -> (a, s)}.
  \pause
\begin{lstlisting}
newtype State s a = State { runState :: s -> (a, s) }
-- krijgt een toestand s, geeft een resultaat a en een nieuwe toestand
\end{lstlisting}
\end{frame}
 
\begin{frame}[fragile]{State is een monad}
\begin{lstlisting}
instance Monad (State s) where
  m >>= f = State $ \s ->
    let (a, s') = runState m s       -- voer m uit
    in  runState (f a) s'            -- kies de volgende stap, geef s' door
instance Applicative (State s) where
  pure x = State $ \s -> (x, s)      -- toestand ongewijzigd
  (<*>)  = ap                        -- ap en liftM: uit Control.Monad
instance Functor (State s) where
  fmap   = liftM
\end{lstlisting}
  \pause
\begin{lstlisting}
get      = State $ \s -> (s, s)      -- get    :: State s s
put s    = State $ \_ -> ((), s)     -- put    :: s -> State s ()
modify f = get >>= put . f           -- modify :: (s -> s) -> State s ()
\end{lstlisting}
\end{frame}
 
\begin{frame}[fragile]{State in actie}
  \begin{columns}[T]
    \column{.4\textwidth}
\begin{lstlisting}
import Control.Monad.State
 
fresh :: State Int Int
fresh = do
  n <- get
  put (n + 1)
  return n
\end{lstlisting}
    \column{.56\textwidth}
\begin{lstlisting}
label :: Tree a -> State Int (Tree (Int, a))
label = traverse (\x -> do n <- fresh
                           return (n, x))
\end{lstlisting}
  \end{columns}
  \pause
\begin{lstlisting}
ghci> evalState (label (Node (Leaf 'a') (Node (Leaf 'b') (Leaf 'c')))) 0
Node (Leaf (0,'a')) (Node (Leaf (1,'b')) (Leaf (2,'c')))
\end{lstlisting}
\end{frame}
 
\begin{frame}[fragile]{State in actie: een deur}
  \begin{columns}[T]
    \column{.5\textwidth}
\begin{lstlisting}
data Door = Open | Closed
  deriving Show
 
step :: String -> State Door ()
step "open"   = put Open
step "close"  = put Closed
step "toggle" = modify toggle
  where toggle Open   = Closed
        toggle Closed = Open
step _        = return ()
\end{lstlisting}
    \column{.48\textwidth}
\begin{lstlisting}
ghci> runState (step "open") Closed
((),Open)
ghci> runState (step "toggle") Open
((),Closed)
\end{lstlisting}
    \pause
\begin{lstlisting}
ghci> execState
        (mapM_ step ["open", "toggle",
                     "toggle"])
        Closed
Open
\end{lstlisting}
  \end{columns}
  \pause
  \vspace{4pt}
  \hs{mapM} is \hs{traverse} (voor monads), \hs{mapM_} gooit de resultaten weg.
\end{frame}

\begin{frame}[fragile]{Samenvatting}
  \centering
  \small
  \begin{tabular}{lll}
    \toprule
    \textbf{Klasse} & \textbf{Kernoperatie} & \textbf{Wat je kan} \\
    \midrule
    \hs|Semigroup|   & \hs|(<>)|      & waarden associatief combineren \\
    \hs|Monoid|      & \hs|mempty|    & \dots\ ook vanuit ``niets'' \\
    \hs|Functor|     & \hs|fmap|      & een gewone functie in een context toepassen \\
    \hs|Applicative| & \hs|pure|, \hs|(<*>)| & onafhankelijke effecten combineren \\
    \hs|Monad|       & \hs|(>>=)|, \hs|join| & afhankelijke effecten na elkaar \\
    \hs|Foldable|    & \hs|foldMap|   & een structuur samenvatten met een monoïde \\
    \hs|Traversable| & \hs|traverse|  & effecten over een structuur, de vorm blijft \\
    \bottomrule
  \end{tabular}

  \vspace{8pt}
  \raggedright
  \begin{alertblock}{Vuistregel}
    Gebruik de \textbf{zwakste} abstractie die doet wat je nodig hebt:
    kan het met \hs|fmap|? Doe dat. Volstaan \hs|pure| en \hs|<*>|? Stop daar.
    Gebruik \hs|>>=| enkel als een stap afhangt van een vorig resultaat. 
  \end{alertblock}
\end{frame}
 
\begin{frame}[standout]
  Vragen?
\end{frame}

\section{Eindquiz}
 
\begin{frame}{Eindquiz: wat gebruik je hier? (1/3)}
  {\small Semigroup, Monoid, Functor, Applicative, Monad, Foldable of Traversable?}
  \begin{enumerate}
    \item elke \code{String} in een \code{Tree String} in hoofdletters zetten
      \uncover<2->{\\ \alert{Functor}: \code{fmap (map toUpper)}. Enkel de waarden veranderen, de vorm blijft.}
    \item foutmeldingen van verschillende checks samenvoegen; geen fouten is ook een geldig resultaat
      \uncover<3->{\\ \alert{Monoid}: \code{[String]} met \code{++}, en \code{mempty = []} betekent ``geen fouten''.}
    \item de hoogste temperatuur uit een \code{NonEmpty Double}
      \uncover<4->{\\ \alert{Semigroup}: \code{Max}. Er is geen zinnig neutraal element voor \code{Double}, maar dat hoeft ook niet.}
    \item nagaan of er ergens in een \code{Tree Int} een negatief getal zit
      \uncover<5->{\\ \alert{Foldable}: \code{any (< 0)}, of \code{getAny . foldMap (Any . (< 0))}. De vorm doet er niet toe.}
  \end{enumerate}
\end{frame}
 
\begin{frame}{Eindquiz: wat gebruik je hier? (2/3)}
  \begin{enumerate}
    \setcounter{enumi}{4}
    \item alle combinaties maken van een maat (S, M, L) en een kleur (rood, blauw)
      \uncover<2->{\\ \alert{Applicative} (lijsten): \code{(,) <\$> maten <*> kleuren}, het cartesisch product.}
    \item elke regel van een bestand parsen naar een \code{Int}; één ongeldige regel en alles faalt
      \uncover<3->{\\ \alert{Traversable}: \code{traverse readMaybe (lines s) :: Maybe [Int]}.}
    \item een gebruiker opzoeken op id, dan zijn afdeling, dan de manager van die afdeling
      (elke stap kan falen)
      \uncover<4->{\\ \alert{Monad}: \code{zoekUser i >\kern0pt>= zoekAfdeling >\kern0pt>= zoekManager}.
        Elke stap heeft het \alert{resultaat} van de vorige nodig.}
    \item een reeks zetten uitvoeren op een spelbord, zonder mutatie
      \uncover<5->{\\ \alert{State} (een Monad) + \code{mapM\_ zet zetten}.}
  \end{enumerate}
\end{frame}
 
\begin{frame}{Eindquiz: wat gebruik je hier? (3/3)}
  \begin{enumerate}
    \setcounter{enumi}{8}
    \item een lijst van transformaties \code{[a -> a]} samenvoegen tot één functie
      \uncover<2->{\\ \alert{Monoid}: functies \code{a -> a} met \code{(.)} en \code{id}, dus \code{foldr (.) id fs}.}
    \item vraag de naam; \alert{enkel} als die \code{"admin"} is, vraag ook een wachtwoord
      \uncover<3->{\\ \alert{Monad}: welke vraag volgt, hangt af van het antwoord. Een Applicative ligt vast op voorhand.}
    \item een formulier valideren en \alert{alle} fouten tonen, niet enkel de eerste
      \uncover<4->{\\ \alert{Applicative} met een \code{Semigroup} voor de fouten. \code{Either} stopt bij de eerste fout,
        en een Monad kan hier niet: \code{>\kern0pt>=} heeft het eerste resultaat nodig om verder te gaan.}
    \item de diepte van een \code{Tree a} berekenen
      \uncover<5->{\\ \alert{Geen van allen}: \code{Foldable} gooit de vorm weg (\code{length} telt bladeren, geen diepte).
        Gewone recursie dus.}
  \end{enumerate}
\end{frame}

\end{document}
